\chapter{Hypernets: Markov Kernels}
\label{ch:qtsl-markov-kernels}

\noindent\draftmark\medskip

\Bookref{ch:qtsl-hypernets}{The hypernet chapter} fixed a
deterministic calculus: ten primitive boxes, one evaluator, one
canonical serialization, one content digest. Stochastic models need
kernels: a diffusion with jumps needs a jump law, a stochastic
Petri-style transition needs a firing law, and a Markov chain simply
\emph{is} its transition kernel. This chapter is about how the
toolkit says ``Markov kernel'' without adding a primitive. A
\emph{Markov-kernel hypernet} is an ordinary well-formed hypernet
together with a designated roster of \emph{noise ports}: inputs read as
fresh uniform draws on $[0,1]$. Its denotation is the Markov kernel
obtained by pushing the uniform reference measure forward through the
deterministic semantics --- every kernel is presented by an explicit
deterministic dilation with its ancilla exposed. Randomness enters at
the interface, never at a box; the wire format, the evaluator, and the
digest scheme never learn that the extension exists. A reader who only
wants to know which named distributions the toolkit provides can go
straight to the coverage dashboard (section~\ref{sec:exp-dist}).

\section{Kernels as dilations}
\label{sec:mk-dilation}

Three consequences of this reading come essentially for free. \emph{Normalization is a
theorem, not a constraint}: every box is total and the reference
measure is a probability, so nothing expressible at this level loses
mass. \emph{Discard is marginalization}: deleting an output wire
marginalizes the kernel. \emph{Copy is sharing, never resampling}:
fan-out of a noise-derived value is an explicit copy row --- both
consumers see one realization --- while an independent resample is a
second noise port, so the two are different graphs with different
digests. The last point is a central subtlety of the probabilistic
string-diagram literature (tossing a coin and copying the outcome is
not tossing two coins~\cite{fritzliang22,bonchi26b}), and the
linearity discipline of the deterministic calculus settles it
structurally, with no new axioms;
Figure~\ref{fig:mk-copy-resample} draws the two graphs.

All three consequences are machine-checked in the toolkit's Lean
tree at the finite-outcome stratum, where a stage is a measurable
outcome map from the noise cube $[0,1]^{N}$ to a finite value set.
Total mass one is the theorem \texttt{sum\_noiseDenotFun}; pushing
the denotation forward along a projection is
\texttt{noiseDenot\_comp}, with discard as the special case that
forgets a coordinate. For
composition, the law runs in the positive direction and in the
negative one: with disjoint noise rosters, a sequential composite
denotes the Kleisli composite of the stage denotations
(\texttt{noiseDenot\_seqComp}), while a shared roster provably
departs from that composite (\texttt{copy\_ne\_kleisli}), so the
disjointness hypothesis is content rather than bookkeeping. All of
these live in the namespace \texttt{SuFunctor.Stochastic},
sorry-free.

\begin{figure}[t]
\centering
\begin{tikzpicture}[x=0.94cm,y=0.82cm,
  box/.style={draw=accentdeep, fill=accent!50, rounded corners=1pt,
    font=\footnotesize\ttfamily, inner sep=2.5pt, minimum height=4.2mm},
  pbox/.style={draw=NoteGray, fill=NoteGray!12, rounded corners=1pt,
    font=\footnotesize\ttfamily, inner sep=2.5pt, minimum height=4.2mm},
  wire/.style={-, line width=0.5pt, draw=black!75},
  port/.style={font=\small\itshape, inner sep=1pt},
  plab/.style={font=\scriptsize, inner sep=1pt, text=black!75}]
\node[port] (u) at (0, 0.75) {$u$};
\node[pbox] (cp) at (1.8, 0.75) {copy$_2$};
\node[box] (f1) at (3.6, 1.5) {$f$};
\node[box] (g1) at (3.6, 0)   {$g$};
\node[port] (o1) at (4.9, 1.5) {$y_1$};
\node[port] (o2) at (4.9, 0)   {$y_2$};
\draw[wire] (u) -- (cp);
\draw[wire] (cp.east) to[out=35,in=180] (f1.west);
\draw[wire] (cp.east) to[out=-35,in=180] (g1.west);
\draw[wire] (f1) -- (o1);
\draw[wire] (g1) -- (o2);
\node[plab, anchor=north] at (2.6, -0.85) {one draw, shared};
\begin{scope}[xshift=7.2cm]
\node[port] (v1) at (0, 1.5) {$u_1$};
\node[port] (v2) at (0, 0)   {$u_2$};
\node[box] (f2) at (1.8, 1.5) {$f$};
\node[box] (g2) at (1.8, 0)   {$g$};
\node[port] (p1) at (3.1, 1.5) {$y_1$};
\node[port] (p2) at (3.1, 0)   {$y_2$};
\draw[wire] (v1) -- (f2);
\draw[wire] (v2) -- (g2);
\draw[wire] (f2) -- (p1);
\draw[wire] (g2) -- (p2);
\node[plab, anchor=north] at (1.55, -0.85) {two independent draws};
\end{scope}
\end{tikzpicture}
\caption{Copy is sharing, never resampling. Fanning a noise port out
through an explicit copy (left) hands both consumers one realization;
an independent redraw is a second noise port (right). The two are
different graphs with different digests, and their denotations
provably differ (\texttt{copy\_ne\_kleisli}).}
\label{fig:mk-copy-resample}
\end{figure}

The evaluator never samples: given
noise values as ordinary rational inputs, the exact evaluator runs
unchanged. Simulation is the evaluator plus an external random-number
driver; certificates consume the kernel's polynomial data and the
reference measure, never samples.

\section{Jump structure: mass, intensity, thinning}
\label{sec:mk-jump}

Continuous time asks for more than normalized kernels. Jump structure
contributes a jump measure and a killing rate to a process's extended
generator~\cite{davis-1993,ethier-kurtz-1986}, and what must be
tracked is \emph{mass}: rates scale it, killing removes it, and
staying put carries it without moving. The kernel vocabulary for
this tier is a small generator language over normalized
kernels --- jump, stay, kill, scale, superposition, and the empty
generator --- in which every term carries a rational
\emph{intensity majorant}, computed by the typing derivation: scale
multiplies it, superposition adds it.
State-dependent intensities $\lambda(x)$ recover from a scalar
majorant $\Lambda$ by thinning: run the clock at rate $\Lambda$ and,
at each ring, branch to the real jump with probability
$\lambda(x)/\Lambda$, staying put otherwise.

Killing is deliberately
distinct from the mass-preserving discard, and stay is deliberately
distinct from zero: a stay contributes nothing to the generator but
carries mass in the grading, which is what thinning and simulation
budgets need to see.

This vocabulary is machine-checked too --- at the finite-state
stratum, in the same sorry-free module family as the dilation
theorems. Every
term elaborates, by structural induction, to a \emph{thinning normal
form} (\texttt{ThinningForm}): a majorant together with one
normalized kernel into the state space extended by two tagged
outcomes, stay and kill. A sub-stochastic mass deficit is therefore
unrepresentable rather than merely discouraged. The elaboration
cannot emit sub-mass, and the declared intensity distributes exactly
over the named channels (\texttt{elaborate\_mass\_partition}, with
\texttt{outflow\_le\_intensity} giving the majorant its meaning).
Elaboration itself preserves the extended generator
(\texttt{generator\_elaborate}). The normal form is deliberately
non-unique, and its two semantics-preserving moves are themselves
theorems: a majorant may be inflated with the slack paid to stay
(\texttt{generator\_dilate}), and self-jump mass may be reclassified
as stay (\texttt{generator\_absorbDiagonal}) --- killing mass may
not, which is the kill/stay separation made mechanical. Composition
is machine-checked at the same stratum (module
\texttt{SuFunctor.Stochastic.JumpAlgebra}): the normal forms carry
the term vocabulary's composite algebra --- superposition,
thinning-scale, coordinate extension, and parallel interleaving, with
the intensity majorants as its grading --- and elaboration is a
homomorphism of that graded algebra, definitionally for the module
operations (\texttt{elaborate\_superpose} and its kin close by
\texttt{rfl}) and in the thinning quotient for interleaving
(\texttt{generator\_elaborate\_interleave}); sequential composition
of jump kernels is no part of this tier's algebra by design ---
successive jumps compose through time, and the Kleisli law for
discrete steps is section~\ref{sec:mk-dilation}'s. The same
tagged form recovers the jump-diffusion generator of
\bookref{ch:qtsl-mathematical-models}{the modelling chapter}. That
chapter's legacy encoding reads a reset mixture's mass deficit as
killing, in Davis's convention; rewriting each deficit into an
explicit kill outcome leaves the legacy transfer increment
unchanged, and the full generator is recovered with nothing
sub-stochastic remaining (\texttt{generator\_eq\_stepMix}). Every
well-typed term thus declares a finite intensity majorant, the same
bounded-intensity condition the non-Zeno arguments elsewhere in this
book compose with.

One consumer is worth naming here: the supervisors of the composition
inventory (section~\ref{sec:cp-supervise}) are ordinary jump-only
Systems --- policy memory is marking, and a policy update's
stochastic content is a firing kernel of this tier --- so control
through enablement adds no kernel vocabulary of its own.

\section{Packaged generators}
\label{sec:mk-packaged}

Two families of packaged generators round out the vocabulary, both in
the packaged-symbol discipline of
\bookref{ch:qtsl-hypernets}{the hypernet chapter} (a definitional
expansion, a projection lemma, and a candidacy rule --- never a new
box).\footnote{Two senses of \emph{generator} meet in this chapter:
the analyst's (section~\ref{sec:mk-jump} presents extended
generators) and the algebraist's (a packaged generator is a named
symbol with a definitional expansion). The word is inherited from
both literatures; the toolkit keeps the senses apart by type.} \emph{Distribution generators} name the classical laws: each is
an alias for a certified piecewise-polynomial quantile table applied
to a noise port, carrying an explicit accuracy budget, so that a model
that draws from a named distribution elaborates into the same hypernet
language every other certificate in this book speaks.
Section~\ref{sec:exp-dist} is their coverage dashboard. The
\emph{stochastic-matrix generator} covers the finite-state case
directly: a kernel between finite-extent wires given by an explicit
tensor of transition probabilities, whose key conformance obligation
is row-stochasticity --- mass one from each source state --- making
finite Markov chains and the jump chains of continuous-time models
first-class citizens rather than encodings.

A status note, in the spirit of the preview note at the front of this
volume. The deterministic substrate is
\bookref{ch:qtsl-hypernets}{the hypernet chapter}'s and is shipped.
The denotational core of the noise-port reading ---
normalization, discard as marginalization, the Kleisli composition
law and its shared-roster failure --- is machine-checked at the
finite-outcome stratum (section~\ref{sec:mk-dilation}), and the jump
vocabulary's semantics --- the intensity grading, the thinning
normal form, its soundness, and the thinning-quotient moves --- are
machine-checked at the finite-state stratum
(section~\ref{sec:mk-jump}). The kernel-language \emph{surface} has
meanwhile settled: the certificate theory
(\path{theories/Qtsl.geolog}) carries it as normative schema. A
kernel citation there names its hypernet by digest together with its
noise roster and three tagged outcome channels: next, stay, and
\emph{cemetery}, the wire format's name for the kill outcome. With
every unit of mass explicitly named by one of the three channels, a
sub-stochastic deficit is unrepresentable, and a pure kill is a
weight-one cemetery branch.

A transition carries a \emph{clock} with
three wings, and the wing decides when the kernel is consulted. A
\emph{racing} clock declares its intensity majorant. At a finite
majorant, its polynomial presentation elaborates definitionally into
the thinning normal form of section~\ref{sec:mk-jump}; the infinite
majorant instead marks the instant-fire boundary case, contributing
no generator term. A \emph{demonic} clock has no timing law
--- whether and when its kernel is consulted belongs to a scheduler
over which every judgment quantifies --- but carries a strictly
positive dwell floor in the system definition: the scheduler may fire
the transition only after that much contiguous dwell inside its guard
region. A row with floor $d$ therefore fires at most $T/d + 1$ times
by horizon $T$ under every resolution, and the demonic wing cannot
produce Zeno behaviour at all. An untimed Petri net transcribes with
unit dwell floors: the transcription is its \emph{discrete-time
unfolding}, embedded in the continuous-time frame. It is faithful
for ordering-and-counting properties, which are all an untimed net
ever asserted. A
\emph{forced} clock fires at exactly its declared positive dwell of
contiguous in-guard time: deterministic delay, completing the triple
--- exponential, immediate, deterministic --- from which
deterministic and stochastic Petri nets are built, so watchdogs,
timeouts, and protocol timers are native rows rather than encodings.
Preemption is by guard reset (whatever fires first wins, and the
forced clock resets with its guard). A same-instant tie between
two forced transitions resolves through the demonic scheduler's
seat: the classical deterministic-net literature needs an ad-hoc
tie-breaking convention exactly where an imprecise-probability
semantics has a seat already waiting. With both dwell floors
strictly positive, the one syntactically visible source of Zeno
behaviour in the vocabulary is an unbounded chain of instant
firings --- precisely what the crossing-budget certificates price.

The stochastic-matrix generator is
carried with row-stochasticity among its named conformance
obligations. The certified quantile tables behind the dashboard's
certified rows exist today as test-suite artefacts, and the dashboard
section states, row by row, what each status rests on and what would
move it. This chapter still owes the settled surface its
presentation: a coming revision states it at the deterministic
calculus's level of rigour, with its soundness against the semantic
universe, and a section on firing kernels alongside.

\input{sections/distributions}
